\documentclass{situationist-preprint}

\title{Hope as a Virtue}
\msid{EXI-15-a}
\msfamily{EXI-15}
\msdiscipline{Meaning, Death and the Absurd}
\mstype{Principal manuscript}
\mskeywords{hope; virtue; despair; presumption; credence; Bovens; Martin; Pettit}

\begin{document}
\maketitle

\begin{abstract}
Is hope a virtue? Its critics, from the Stoics to Nietzsche, regard it as a weakness that binds us to what we cannot control and tempts us to self-deception; its defenders, from Aquinas to recent moral psychology, regard it as an excellence of agents who must act under uncertainty. We argue that hope is a virtue when it is proportioned to probability without collapsing into it. Following Bovens, Martin and Pettit, we analyze hope as desire for an outcome, a credence strictly between zero and one, and a practical stance that treats the outcome's possibility as a reason for planning, attention and effort. The credence is fixed by the evidence and is not chosen; the stance is under the agent's control, as attention is. The virtue lies in the relation between them: the stance must be sensitive to the probability, so that hope for the nearly impossible does not drive large investments, but it must not be fixed by the probability alone, since two agents with the same credence may rightly hope or not according to what is at stake and what they can do. Despair and presumption are the vices on either side.
\end{abstract}

\section{The problem}

For the Stoics, hope and fear are twin disturbances: both attach us to future things not in our power, and the wise person is free of both. Nietzsche read the hope left in Pandora's jar as the worst of evils, since it prolongs torment. Aquinas, by contrast, counted hope among the theological virtues and identified two vices opposed to it: despair, which gives up a good that can be attained, and presumption, which counts on attaining it without regard to the conditions of its attainment (\emph{Summa Theologiae} II-II, qq.~17--21). The question of this family is whether hope, in its ordinary, secular form, is a virtue in the sense of a disposition that makes its bearer act and feel well.

The answer depends on what hope is. On the \emph{standard account}, to hope that $p$ is to desire that $p$ and to assign $p$ a probability between zero and one \citep{Downie1963}. The standard account cannot distinguish two people with the same desire and the same credence, one of whom hopes and the other despairs. \citet{Bovens1999} added a mental imaging of what it would be like if $p$; \citet{Pettit2004} added a resolve to act as if $p$ will obtain, or at least as if it might, which he called substantial hope; \citet{Martin2014} proposed an incorporation analysis, on which the hoper treats the probability of the outcome as licensing her to treat her desire as a reason for planning and feeling. We build on these analyses.

\section{Framework}

\begin{definition}[Hope]\label{d:hope}
An agent \emph{hopes} that $p$ if (i) she desires that $p$; (ii) her credence in $p$ is strictly between zero and one; and (iii) she takes up a \emph{hopeful stance} toward $p$: she treats the possibility of $p$ as a reason to plan for it, to attend to the ways in which it might come about, and to make efforts toward it where efforts are possible.
\end{definition}

\begin{principle}[Credence is evidential]\label{p:cred}
The rational credence in $p$ is fixed by the agent's evidence (EPI-17), and the agent does not control her credence directly (EPI-11).
\end{principle}

\begin{principle}[The stance is controlled]\label{p:stance}
The hopeful stance is under the agent's control, since it consists in what she attends to, plans for and attempts, and attention is under voluntary control (EPI-11).
\end{principle}

\begin{definition}[Virtue]\label{d:virtue}
A \emph{virtue} is a stable disposition to feel and act in the right way, in the right circumstances and to the right degree, flanked by vices of deficiency and excess (Aristotle, \emph{Nicomachean Ethics} II.6).
\end{definition}

Principles~\ref{p:cred} and~\ref{p:stance} together fix the location of hope's rationality. Hope does not license a raised credence: a hopeful patient who comes to believe her recovery is likely because she wants it has not hoped well but believed badly. What hope adds to credence is a practical commitment, in the way that, on the view of EPI-23, James's right to believe licenses not a different credence but a commitment to act on a hypothesis. The objects of evaluation are thus the stance and its relation to the credence.

\section{The argument}

\begin{definition}[Proportioned hope]\label{d:prop}
Hope that $p$ is \emph{proportioned to probability without collapsing into it} if (a) the strength of the hopeful stance, measured by the resources of attention, planning and effort it commits, increases with the credence in $p$, other things equal; and (b) the strength is not a function of the credence alone, but depends also on the value of $p$ and on the agent's capacity to make $p$ more likely.
\end{definition}

\begin{lemma}\label{l:sensitive}
Hope whose stance is insensitive to probability is a vice of excess.
\end{lemma}

\begin{proof}
If the stance is insensitive to probability, an agent invests as much attention, planning and effort in a remote possibility as in a likely one with the same value. Resources spent on a hope are unavailable for other ends, so she systematically misallocates them toward what is unlikely to be realized. This is Aquinas's presumption, in its secular form: counting on an outcome without regard to the conditions of its attainment. It is acting to the wrong degree (Definition~\ref{d:virtue}).
\end{proof}

\begin{lemma}\label{l:collapse}
Hope whose stance is fixed by probability alone is deficient.
\end{lemma}

\begin{proof}
Suppose the stance were a function of the credence alone. Then two agents with the same credence in $p$ would be required to take the same stance whatever the value of $p$ and whatever they could do about it. But a parent whose child has a one-in-five chance of recovery has reason to sustain treatment, attention and planning for recovery that a spectator with the same credence about a horse race lacks, because recovery is worth more and because her efforts can affect the outcome. A disposition that ignores these differences gives up goods that can be attained where they are valuable and her agency counts. This is despair in its secular form, and it is acting to the wrong degree in the other direction.
\end{proof}

\begin{theorem}\label{t:main}
Hope is a virtue when it is proportioned to probability without collapsing into it.
\end{theorem}

\begin{proof}
By Lemmas~\ref{l:sensitive} and~\ref{l:collapse}, hope fails as a disposition when its stance is insensitive to probability (excess) and when it is fixed by probability alone (deficiency). Hope proportioned in the sense of Definition~\ref{d:prop} avoids both: it varies with the evidence, and it also varies with value and agency, which are the other things that determine the right degree of investment in an uncertain outcome. A stable disposition to hope in this way is a disposition to act and feel to the right degree in the right circumstances, between two vices; by Definition~\ref{d:virtue}, it is a virtue. Since the credence component is fixed by the evidence (Principle~\ref{p:cred}), the disposition does not require any departure from epistemic rationality; it consists in the governance of the stance, which is within the agent's control (Principle~\ref{p:stance}).
\end{proof}

\section{Objections and replies}

\begin{objection}[Stoic]\label{o:stoic}
Hope attaches us to outcomes beyond our control, and so exposes us to disappointment and anxiety. The wise person accepts what comes.
\end{objection}

\begin{reply}
The Stoic objection targets hope whose stance outruns the agent's agency and the evidence, which Lemma~\ref{l:sensitive} already counts as a vice. Proportioned hope invests in an outcome in proportion to its probability and to what one can do about it. It exposes the agent to disappointment, but so does every attachment to an uncertain good, including love and the pursuit of achievement, and the Stoic remedy of detachment from all of them gives up too much (see EXI-07 on the dependence of happiness on caring about things other than one's own state).
\end{reply}

\begin{objection}[Self-deception]\label{o:deceive}
Hope sustains itself by turning attention away from evidence against the hoped-for outcome. Since belief follows attention (EPI-11), hope will corrupt credence.
\end{objection}

\begin{reply}
The objection identifies a real danger, and it is the reason the virtue requires proportion. Attention directed to the ways in which an outcome might come about is not attention withdrawn from the evidence about how likely it is. The virtuous hoper plans for recovery while reading the prognosis honestly, which is the combination that \citet{McGeer2004} calls good hope and distinguishes from wishful hope. Where the stance corrupts the credence, hope has failed in exactly the respect Definition~\ref{d:prop}(a) requires.
\end{reply}

\begin{objection}[Hope against hope]\label{o:against}
We admire people who hope when the odds are very poor: the prisoner who survives by hoping for release, the activist who works for a cause that seems lost. Proportion to probability condemns them.
\end{objection}

\begin{reply}
Definition~\ref{d:prop} does not forbid hope at low probability; it forbids a stance insensitive to probability. Where the outcome is of great value and one's efforts are among its causes, clause (b) supports substantial hope at low credence, which is why the prisoner's and the activist's hope can be virtuous. Pettit's point that hope can sustain agency when evidence offers no handhold \citep{Pettit2004} is a point about this clause. What the theorem condemns is the stance that would be the same whatever the odds.
\end{reply}

\section{Conclusion}

Hope is desire plus a credence short of certainty plus a practical stance, and its virtue lies in the stance. The credence is set by the evidence and should not be bent by the desire. The stance should rise and fall with the probability, but also with what is at stake and with what one can do. Hope that ignores probability is presumption, and hope that is nothing but probability gives up goods that agency could reach; between them lies a virtue.

\begin{thebibliography}{99}\small
\bibitem[Bovens(1999)]{Bovens1999} Bovens, L. (1999). The value of hope. \emph{Philosophy and Phenomenological Research}, 59(3), 667--681.
\bibitem[Downie(1963)]{Downie1963} Downie, R.~S. (1963). Hope. \emph{Philosophy and Phenomenological Research}, 24(2), 248--251.
\bibitem[Martin(2014)]{Martin2014} Martin, A.~M. (2014). \emph{How We Hope: A Moral Psychology}. Princeton University Press.
\bibitem[McGeer(2004)]{McGeer2004} McGeer, V. (2004). The art of good hope. \emph{The Annals of the American Academy of Political and Social Science}, 592, 100--127.
\bibitem[Pettit(2004)]{Pettit2004} Pettit, P. (2004). Hope and its place in mind. \emph{The Annals of the American Academy of Political and Social Science}, 592, 152--165.
\end{thebibliography}

\end{document}
